\documentclass[../main.tex]{subfiles}
\begin{document}

\section{Summary}

This design project has addressed the critical optimization challenges in Multilingual Text-Based Person Search (MTBPS), specifically focusing on the low-resource context of the Vietnamese language. Through a rigorous analysis of the gradient dynamics of the standard Sigmoid loss (N-ITC), we identified a "vanishing gradient" problem for semi-hard negative samples, which hinders the model's discriminative power.

\noindent To resolve this, we proposed a novel Hard Negative-Aware Optimization Framework that integrates Cross-modal Circle Loss with Low-Rank Adaptation (LoRA). Our key contributions and findings are summarized as follows:
\begin{itemize}
    \item Geometric Refinement: The integration of Circle Loss successfully reshaped the embedding space, enforcing a spherical geometry with clear margins between identities, as evidenced by the geometric visualization.

\item Computational Efficiency: By utilizing LoRA, we reduced the memory footprint significantly, enabling a 3x increase in batch size (from 8 to 24) on consumer hardware. This was proven to be a prerequisite for effective contrastive learning.

\item State-of-the-Art Performance: The proposed curriculum-based strategy achieved a new SOTA Rank-1 accuracy of 51.30\% on the VnPersonSearch benchmark and 71.85\% on CUHK-PEDES, validating the robustness of the method across both low and high-resource settings.
\end{itemize}

\section{Limitations}
Despite the promising results, the current framework still faces certain technical limitations that constrain its potential performance:

Fixed Aspect Ratio (Square Input): Currently, all input images are resized to a fixed resolution of 256×256 to match the mSigLIP pre-training protocol. However, pedestrian images are naturally rectangular (tall and narrow). Forcing them into a square aspect ratio distorts the human body shape and aspect ratio, potentially losing critical morphological details required for re-identification.

Limited Feature Plasticity with LoRA: While LoRA is memory-efficient, it restricts the optimization to a low-rank subspace (r=32). This means the model cannot drastically alter its deep feature representations to adapt to extremely difficult domain shifts or highly idiosyncratic Vietnamese attributes. It represents a trade-off between efficiency and maximum theoretical plasticity.

Suboptimal Batch Size: Although LoRA allowed us to increase the batch size to 24, this number is still relatively small for Contrastive Learning, which typically benefits from batch sizes in the hundreds or thousands (e.g., CLIP uses 32k). A small batch size limits the number of negative samples available in each iteration, reducing the probability of encountering "true" hard negatives during training.

\section{Suggestion for Future Works }

Based on the identified limitations and the results of this design project, the future development direction for the Bachelor’s Graduation Thesis will focus on the following key areas:
\begin{itemize}
\item Adaptive Resolution Strategy: We plan to implement a dynamic resolution mechanism or a "patch-keeping" strategy that preserves the original aspect ratio of pedestrian images. This ensures that fine-grained details (e.g., shoes, tall hairstyles) are not distorted by resizing artifacts.

\item Scaling Batch Size via Gradient Accumulation: To overcome the GPU memory constraint while maintaining the benefits of large-batch contrastive learning, we will implement Gradient Cache or Gradient Accumulationtechniques. This will allow us to simulate a much larger effective batch size (e.g., 64 or 128), providing a richer pool of negative samples for the Circle Loss to mine.

\item Exploring Advanced PEFT Methods: Beyond LoRA, we will investigate other Parameter-Efficient Fine-Tuning methods such as DoRA (Weight-Decomposed Low-Rank Adaptation) or AdapterFusion. These methods may offer a better balance between parameter efficiency and the model's ability to learn complex, non-linear feature transformations for the Vietnamese language.

\item Extending to Video-Based Retrieval: Given the temporal nature of surveillance data, we aim to extend the framework to support Video-to-Text retrieval, leveraging temporal attention mechanisms to aggregate cues across multiple frames. 
\end{itemize}
\end{document}